% !PRINTED 50
\weberpage{50}

Wenn in der Formel (1) $\varphi$ in $-\varphi$ verwandelt wird, so geht $u$ in $-u$ über, es ist also $\mathrm{am}\,(-u) = -\mathrm{am}\,u$ und folglich
\[
\mathrm{sn}\,(-u) = -\mathrm{sn}\,u,\quad \mathrm{cn}\,(-u) = \mathrm{cn}\,u,\quad \mathrm{dn}\,(-u) = \mathrm{dn}\,u,
\]
d. h. es ist sn $u$ eine ungerade, cn $u$, dn $u$ sin gerade Funktionen. % [sic] \textquotedbl{}sin\textquotedbl{} printed for \textquotedbl{}sind\textquotedbl{}
Setzt man also
\[
\frac{d\varphi_2}{\Delta\varphi_2} = du,\quad \frac{d\varphi_3}{\Delta\varphi_3} = dv,\quad \frac{d\varphi_1}{\Delta\varphi_1} = -d(u+v),
\]
so lassen sich die Formeln (17), §\,13 anwenden, und man erhält die \gesperrt{Additionstheoreme}:
\weq{6}{\begin{gathered}
\mathrm{sn}\,(u+v) = \dfrac{\mathrm{sn}\,u\,\mathrm{cn}\,v\,\mathrm{dn}\,v + \mathrm{sn}\,v\,\mathrm{cn}\,u\,\mathrm{dn}\,u}{1 - \varkappa^2\mathrm{sn}^2u\,\mathrm{sn}^2v},\\
\mathrm{cn}\,(u+v) = \dfrac{\mathrm{cn}\,u\,\mathrm{cn}\,v - \mathrm{sn}\,u\,\mathrm{sn}\,v\,\mathrm{dn}\,u\,\mathrm{dn}\,v}{1 - \varkappa^2\mathrm{sn}^2u\,\mathrm{sn}^2v},\\
\mathrm{dn}\,(u+v) = \dfrac{\mathrm{dn}\,u\,\mathrm{dn}\,v - \varkappa^2\mathrm{sn}\,u\,\mathrm{sn}\,v\,\mathrm{cn}\,u\,\mathrm{cn}\,v}{1 - \varkappa^2\mathrm{sn}^2u\,\mathrm{sn}^2v}.
\end{gathered}}
Man kann daraus die drei elliptischen Funktionen für beliebige Werte des Arguments \gesperrt{eindeutig} berechnen, wenn sie für irgend ein noch so kleines Gebiet um den Nullpunkt bestimmt sind. Diesen Satz nimmt \gesperrt{Weierstrass} zum Ausgangspunkt der Theorie der elliptischen Funktionen.

Setzt man
\[
\int_0^{\frac{\pi}{2}} \frac{d\varphi}{\Delta\varphi} = K,
\]
so ist $\mathrm{sn}\,K = 1$, $\mathrm{cn}\,K = 0$, $\mathrm{dn}\,K = \sqrt{1-\varkappa^2} = \varkappa'$, und es ergeben also die Formeln (6), indem man $v = K$ setzt:
\weq{7}{\begin{gathered}
\mathrm{sn}\,(u+K) = \dfrac{\mathrm{cn}\,u}{\mathrm{dn}\,u},\\
\mathrm{cn}\,(u+K) = \dfrac{-\varkappa'\,\mathrm{sn}\,u}{\mathrm{dn}\,u},\\
\mathrm{dn}\,(u+K) = \dfrac{\varkappa'}{\mathrm{dn}\,u},
\end{gathered}}
und wenn man diese Formel noch einmal anwendet:
\weq{8}{\begin{gathered}
\mathrm{sn}\,(u+2K) = -\mathrm{sn}\,u,\\
\mathrm{cn}\,(u+2K) = -\mathrm{cn}\,u,\\
\mathrm{dn}\,(u+2K) = \mathrm{dn}\,u,
\end{gathered}}
%CONT>>
